\documentclass{article}
\usepackage[utf8]{inputenc}
\usepackage{multirow}
\begin{document}

Kosmické rychlosti:
\begin{enumerate}
    \item kruhová rychlost, těsně nad povrchem Země $\doteq 7.9\,\mathrm{km}\cdot\mathrm{s}^{-1}$
        \begin{equation}
            v_k = \sqrt{\frac{GM}{r}} ~,
        \end{equation}
    \item úniková rychlost, z povrchu Země $\doteq 11.2\,\mathrm{km}\cdot\mathrm{s}^{-1}$
        \begin{equation}
            v_p = \sqrt{2}\cdot v_k ~,
        \end{equation}
\end{enumerate}

Zajímavé názvy apsid (pericenter a apocenter, nebo periapsid a apoapsid) z české wiki:\\
\begin{tabular}{lll}
Země                               & perigeum (přízemí)      & apogeum (odzemí) \\
\hline
\multirow{3}{.2\linewidth}{Měsíc}  & periselenium            & aposelenium \\
                                   & pericynthion            & apocynthion \\
                                   & perilun(ium)            & apolun(ium) \\
\hline
Slunce                             & perihel(ium) (přísluní) & afel(ium) (odsluní) \\
\hline
Venuše                             & pericytherion           & apocytherion \\
\hline
\multirow{2}{.2\linewidth}{Mars}   & periareon               & apareon \\
                                   & periareum               & apoareum \\
\hline
Jupiter                            & perijov(ium)            & apojov(ium) \\
\hline
\multirow{2}{.2\linewidth}{Saturn} & perikron(ium)           & apokron(ium) \\
                                   & perisaturnium           & aposaturnium \\
\hline
hvězda                             & periastron              & apoastron \\
\hline
střed Galaxie                      & perigalaktikum          & apogalaktikum \\
\hline
černá díra                         & peribothon              & apobothon
\end{tabular}

Apsidy v problému dvou těles pevné! Ve Sluneční soustavě:
\begin{itemize}
    \item apsidy se obecně stáčejí díky interakci s dalšími tělesy,
    \item podle odchylek trajektorie Uranu určil Urbain Le Verrier polohu neznámé planety,
    \item 23. září 1846 z berlínské observatoře nalezen Neptun s odchylkou pod 1 stupeň, (pozoroval ho už Galileo 1612 a 1613)
    \item u Merkuru celkem 5600 úhlových vteřin za století (cca 400 orbit),
    \item z toho 43 úhlové vteřiny za století plynou z obecné relativity
\end{itemize}

\end{document}
